\subsubsection{Mathematical Relationship Between Efficiency Gap and Proportionality}
\label{sec:eg-proportionality-connection}

Before presenting empirical proportionality findings, we clarify the mathematical relationship between efficiency gap and proportional representation. These concepts are related but distinct, and conflating them leads to analytical confusion.

\textbf{Definitions:}
\begin{itemize}
    \item \textbf{Efficiency gap (EG)}: Measures \emph{partisan asymmetry} in wasted votes
    \item \textbf{Proportionality gap}: Measures \emph{deviation from proportional representation} (seat share - vote share)
\end{itemize}

\textbf{Approximate relationship:}
For competitive elections (vote share near 50\%), efficiency gap approximates twice the proportionality gap:

\begin{equation}
\text{EG} \approx 2 \times (\text{Seat Share} - \text{Vote Share})
\end{equation}

\textbf{Derivation (simplified):}
Consider a state with $V_D$ Democratic votes and $V_R$ Republican votes, producing $S_D$ Democratic seats and $S_R$ Republican seats.

\begin{align*}
\text{EG} &= \frac{\text{Wasted}_R - \text{Wasted}_D}{\text{Total Votes}} \\
&\approx \frac{(V_R - S_R \cdot V_{\text{threshold}}) - (V_D - S_D \cdot V_{\text{threshold}})}{V_D + V_R}
\end{align*}

Where $V_{\text{threshold}} = (V_D + V_R) / (S_D + S_R) \cdot 0.5$ is the average votes needed to win a district.

With algebraic manipulation (see Stephanopoulos \& McGhee, 2015, Appendix B):

\begin{equation}
\text{EG} \approx 2 \times \left(\frac{S_D}{S_D + S_R} - \frac{V_D}{V_D + V_R}\right)
\end{equation}

\textbf{Interpretation:} Efficiency gap captures proportionality deviations with factor-of-two scaling. A $-3.2\%$ efficiency gap implies $\approx -1.6$ percentage point proportionality gap (Democrats win $\approx$1.6\% more seats than their vote share predicts).

\textbf{Example calculation for our findings:}

\begin{align*}
\text{Algorithmic plans:} \quad & \text{EG} = -3.2\% \\
& \Rightarrow \text{Proportionality gap} \approx -1.6\% \\
& \text{If Dem vote share = 52\%, seat share} \approx 52\% + 1.6\% = 53.6\%
\end{align*}

Our proportionality analysis (Section~\ref{sec:proportionality-analysis}) reports Democrats winning 54.1\% seats with 51.3\% votes, yielding $+2.8$ pp proportionality gap. The EG-based prediction ($-3.2\% / 2 = -1.6\%$ proportionality gap, but applied to 51.3\% vote share) gives:
\begin{equation}
51.3\% + 1.6\% = 52.9\% \approx 53-54\% \text{ (close to observed 54.1\%)}
\end{equation}

The approximation holds within $\approx$1 percentage point, confirming the EG-proportionality relationship.

\textbf{Why the approximation breaks down:}

The EG $\approx$ 2$\times$(seat share - vote share) relationship assumes:
\begin{enumerate}
    \item \textbf{Vote shares near 50\%}: Works well for competitive states (45-55\% vote share); breaks down for landslides
    \item \textbf{Uniform swing}: Assumes proportional vote changes across districts; violated when specific regions swing differently
    \item \textbf{Similar turnout}: Assumes equal turnout across districts; urban-rural turnout differentials introduce error
\end{enumerate}

For our analysis of 15 competitive states (Democratic vote shares 48-54\%), the approximation is reasonably accurate. For less competitive states, efficiency gap and proportionality gap can diverge.

\textbf{Conceptual distinction:}

Despite mathematical relationship, efficiency gap and proportionality measure different fairness concepts:

\begin{itemize}
    \item \textbf{Efficiency gap}: Measures \emph{partisan symmetry}. Question: ``If parties' vote shares were reversed, would seat shares reverse equally?''
    \item \textbf{Proportionality}: Measures \emph{majoritarian representation}. Question: ``Does the party winning 52\% votes win 52\% seats?''
\end{itemize}

\textbf{Example where they diverge:}

Consider a state where:
\begin{itemize}
    \item Democrats win 52\% votes, 48\% seats (proportionality gap: $-4\%$)
    \item Republicans would win 48\% seats with 52\% votes (by symmetry)
\end{itemize}

This plan has perfect partisan symmetry (zero bias) but imperfect proportionality ($-4\%$ gap for both parties). Conversely, a plan can be perfectly proportional (52\% votes $\rightarrow$ 52\% seats) but show partisan bias if 48\% votes for the other party yields only 45\% seats.

\textbf{Implications for our findings:}

Our $-3.2\%$ efficiency gap for algorithmic plans indicates:
\begin{enumerate}
    \item \textbf{Proportionality implication}: Democrats win $\approx$1.6 percentage points more seats than proportionality predicts
    \item \textbf{Symmetry implication}: Democrats have slight advantage; Republicans face symmetric disadvantage
    \item \textbf{Geographic cause}: Urban Democratic concentration creates this asymmetry under compact districting
\end{enumerate}

The positive efficiency gap ($+5.1\%$) for enacted plans indicates:
\begin{enumerate}
    \item \textbf{Proportionality implication}: Republicans win $\approx$2.5 percentage points more seats than proportionality predicts
    \item \textbf{Symmetry implication}: Republicans advantaged; Democrats disadvantaged
    \item \textbf{Mixed causes}: Geography ($-3.2\%$ baseline) plus manipulation ($+8.3\%$ deviation from baseline)
\end{enumerate}

\textbf{Connecting to mean-median difference:}

Mean-median difference provides another proportionality-related metric:

\begin{equation}
\text{MMD} = \text{Median}(\text{Party vote share}) - \text{Mean}(\text{Party vote share})
\end{equation}

When Democrats' mean vote share exceeds median, they are packed into high-margin districts while losing many close races---a packing/cracking pattern. MMD connects to wasted votes (efficiency gap) as follows:

\begin{itemize}
    \item \textbf{High positive MMD}: Party packed into landslide districts $\rightarrow$ surplus wasted votes $\rightarrow$ negative EG
    \item \textbf{Low or negative MMD}: Party wins many close districts $\rightarrow$ minimal wasted votes $\rightarrow$ positive EG
\end{itemize}

Our findings (Section~\ref{sec:proportionality-analysis}):
\begin{itemize}
    \item Algorithmic plans: MMD = $+0.8$ pp $\rightarrow$ modest Democratic packing $\rightarrow$ $-3.2\%$ EG
    \item Enacted plans: MMD = $+4.1$ pp $\rightarrow$ severe Democratic packing $\rightarrow$ (should produce negative EG, but enacted plans show $+5.1\%$ EG due to additional Republican cracking)
\end{itemize}

The discrepancy between MMD (suggests Democratic disadvantage from packing) and EG (shows Republican advantage) in enacted plans reveals that enacted maps use \emph{both} packing (high MMD) \emph{and} cracking (shifts EG toward Republicans). Algorithmic plans show only packing from geographic concentration.

\textbf{Summary:}

The efficiency gap framework provides theoretical foundation for our proportionality analysis. The approximate relationship (EG $\approx$ 2$\times$ proportionality gap) allows translating efficiency gap findings into seat-share predictions:

\begin{itemize}
    \item Algorithmic $-3.2\%$ EG predicts Democrats win $\approx$54\% seats with 52\% votes
    \item Enacted $+5.1\%$ EG predicts Republicans win $\approx$52.5\%$ seats with 48\% votes
\end{itemize}

These predictions align with empirical proportionality findings (Section~\ref{sec:proportionality-analysis}), confirming the robustness of efficiency gap as a partisan fairness measure when contextualized appropriately. The connection also clarifies that our $-3.2\%$ algorithmic baseline is not ``partisan neutrality'' (which would be 0\% EG and perfect proportionality) but rather ``geographic neutrality'' (the minimum bias achievable under compact districting given residential sorting patterns).
